\subsection{一般结果}

设 X 为局部紧空间，局部紧群 H 在其上通过 \((x,\xi)\mapsto x\xi\)（\(x\in X\)，\(\xi\in H\)）连续且适当地右作用。由 H 在 X 中定义的等价关系是开的（GT, III, §2, No.~4, 引理 2），且 X/H 是豪斯多夫空间（同上，§4, No.~2, 命题 3），因而是局部紧的（GT, I, §10, No.~4, 命题 10）。记 \(\pi\) 为从 X 到 X/H 的典范映射。X 的子集 Y 的饱和集为 \(\mathrm{YH}=\pi^{-1}\big(\pi(\mathrm{Y})\big)\)。若 K 是 X 的紧子集，则 \(\pi(\mathrm{K})\) 是紧的，而 K 的饱和集 \(\pi^{-1}\big(\pi(\mathrm{K})\big)\) 在 X 中是闭的。X/H 的每个紧子集都是 X 中某个紧子集在 \(\pi\) 下的像（GT, I, §10, No.~4, 命题 10）。始终固定 H 上的左哈尔测度 \(\beta\)。

设 \(\chi\) 为 \(\mathbf{H}\) 在 \(\mathbf{R}_+^*\) 中的连续表示。若函数 \(g\) 定义在 \(\mathbf{X}\) 上并满足 \(g(x\xi) = \chi (\xi)g(x)\) 对所有 \(x\in \mathbf{X}\) 和 \(\xi \in \mathbf{H}\) 成立，则其支集 \(\mathbf{S}\) 在 \(\mathbf{H}\) 下不变，因而可写为 \(\pi^{-1} (\pi (\mathrm{S}))\)。我们用 \(\mathcal{K}^{\chi}(\mathbf{X})\) 表示由连续实值函数 \(g\) 组成的里斯空间，这些函数定义在 \(\mathbf{X}\) 上，满足 \(g(x\xi) = \chi (\xi)g(x)\)（\(x\in \mathbf{X},\xi \in \mathbf{H}\)），且其支集是 \(\mathbf{X}\) 中某个紧子集的饱和集；用 \(\mathcal{K}_+^\chi (\mathbf{X})\) 表示 \(\geqslant 0\) 的 \(\mathcal{K}^{\chi}(\mathbf{X})\) 元素集合。特别地，\(\mathcal{K}^1 (\mathbf{X})\) 正是定义在 \(\mathbf{X}\) 上、在轨道上为常值且支集是某个紧子集的饱和集的连续函数集合。

命题 1. --- 设 f 是 X 上的连续实值函数，其支集 S 与 X 的每个紧子集的饱和集的交都是紧的。

\begin{enumerate}
\def\labelenumi{\alph{enumi})}
\tightlist
\item
  对于每个 \(x \in \mathbf{X}\)，函数 \(\xi \mapsto f(x\xi)\) 在 \(\mathbf{H}\) 上属于 \(\mathcal{K}(\mathbf{H})\)；令
\end{enumerate}

\[
f ^ {\chi} (x) = \int_ {\mathrm{H}} f (x \xi) \chi (\xi) ^ {- 1} d \beta (\xi).\tag{1}
\]

\begin{enumerate}
\def\labelenumi{\alph{enumi})}
\setcounter{enumi}{1}
\item
  函数 \(f^{\chi}\) 是连续的，在 SH 外为零，并满足 \(f^{\chi}(x\xi) = \chi(\xi)f^{\chi}(x)\)。
\item
  若 \(g\) 是 \(\mathbf{X}\) 上的连续实值函数且满足 \(g(x\xi) = \chi(\xi)g(x)\)，则 \((fg)^{\chi} = f^{1}g\)（其中 \(f^{1}\) 由公式（1）给出，公式中的 \(\chi\) 换成表示 \(\xi \mapsto 1\)，该表示是 \(\mathbf{H}\) 在 \(\mathbf{R}_{+}^{*}\) 中的表示）。
\item
  若 \(\eta \in \mathbf{H}\)，则 \(\left(\pmb{\delta}(\eta)f\right)^{\chi} = \chi (\eta)\Delta_{\mathrm{H}}(\eta)^{-1}f^{\chi}\)。
\end{enumerate}

令 \(x_0 \in X\)，并令 \(V\) 为 \(x_0\) 在 \(X\) 中的紧邻域。集合 \(\xi \in H\) 中满足 \(V\xi\) 与 \(S\) 相交的元素，同时也是集合 \(\xi \in H\) 中满足 \(V\xi\) 与 \(S \cap VH\) 相交的元素，因而在 \(H\) 中是紧的，因为 \(S \cap VH\) 是紧的且 \(H\) 在 \(X\) 上适当地作用（GT, III, §4, No.~5, 定理 1）；于是第 1 节第 1 号的引理 1 证明了 a) 以及 \(f^{\chi}\) 的连续性。b) 的其余部分显然。最后，c) 和 d) 由以下计算得到：

\[
\begin{array}{c} (f g) ^ {\chi} (x) = \int_ {\mathrm{H}} f (x \xi) g (x \xi) \chi (\xi) ^ {- 1} d \beta (\xi) = \int_ {\mathrm{H}} f (x \xi) g (x) \chi (\xi) \chi (\xi) ^ {- 1} d \beta (\xi) \\ = g (x) \int_ {\mathrm{H}} f (x \xi) d \beta (\xi) = g (x) f ^ {1} (x) \end{array}
\]

\[
\begin{array}{r l} & {\left(\delta (\eta) f\right) ^ {\chi} (x) = \int_ {\mathrm{H}} f (x \xi \eta) \chi (\xi) ^ {- 1} d \beta (\xi)} \\ & {\qquad = \Delta_ {\mathrm{H}} (\eta) ^ {- 1} \int_ {\mathrm{H}} f (x \xi) \chi (\xi \eta^ {- 1}) ^ {- 1} d \beta (\xi)} \\ & {\qquad = \chi (\eta) \Delta_ {\mathrm{H}} (\eta) ^ {- 1} \int_ {\mathrm{H}} f (x \xi) \chi (\xi) ^ {- 1} d \beta (\xi).} \end{array}
\]

命题 2. --- 映射 \(f \mapsto f^{\chi}\) 从 \(\mathcal{K}(X)\) 到 \(\mathcal{K}^{\chi}(X)\) 是线性的，且 \(\mathcal{K}(X)\)（分别为 \(\mathcal{K}_{+}(X)\)）的像是 \(\mathcal{K}^{\chi}(X)\)（分别为 \(\mathcal{K}_{+}^{\chi}(X)\)）。

线性性显然。显然 \(f^{\chi} \geqslant 0\) 对 \(f \geqslant 0\) 成立。因此只需应用以下引理：

引理 1. --- 设 K 为 X 的紧子集，u 为 \(\mathcal{K}_{+}(X)\) 中满足 \(u(x) > 0\) 对 \(x \in K\) 成立的函数。设 \(g \in \mathcal{K}^{\chi}(X)\) 满足 \(\operatorname{Supp} g \subset KH\)。

\begin{enumerate}
\def\labelenumi{\alph{enumi})}
\item
有 \(\inf_{x\in \mathrm{KH}}u^1 (x) > 0\)
\item
函数 \(h\) 在 KH 上等于 \(g / u^1\)，在 X - KH 上等于 0，并属于 \(\mathcal{K}^{\chi}(X)\)。
\item
  \(g = (uh)^{\chi}\)。
\end{enumerate}

有 \(u^1(x) > 0\) 对 \(x \in \mathbf{K}\) 成立，因而 \(\inf_{x \in \mathrm{KH}} u^1(x) = \inf_{x \in \mathbf{K}} u^1(x) > 0\)。断言 b) 立即由此得到。最后，根据命题 1 c)，\((uh)^{\chi} = u^{1}h\)，且显然 \(u^1h = g\)。

令 I 为定义在 \(\mathcal{K}^{\chi}(\mathbf{X})\) 上的相对有界线性形式（第二章，§2，第 2 号）。于是 \(f \mapsto \mathrm{I}(f^{\chi})\) 是定义在 \(\mathcal{K}(\mathbf{X})\) 上的相对有界线性形式，即测度 \(\mu_{\mathrm{I}}\) 定义在 \(\mathbf{X}\) 上。映射 \(\mathrm{I} \mapsto \mu_{\mathrm{I}}\) 由命题 2 知是单射。由此得到的测度 \(\mu_{\mathrm{I}}\) 定义在 \(\mathbf{X}\) 上，可如下刻画：

命题 3. --- 设 \(\mu\) 是 X 上的测度。以下条件等价：

\begin{enumerate}
\def\labelenumi{\alph{enumi})}
\item
  存在定义在 \(\mathcal{K}^{\chi}(\mathbf{X})\) 上的相对有界线性形式 I，使得 \(\mathrm{I}(f^{\chi}) = \mu(f)\) 对所有 \(f\in \mathcal{K}(\mathbf{X})\) 成立。
\item
  \(\delta (\xi)\mu = \chi (\xi)^{-1}\Delta_{\mathrm{H}}(\xi)\mu\) 对所有 \(\xi \in \mathbf{H}\) 成立。
\item
  对所有 f、g 属于 \(\mathcal{K}(\mathbf{X})\)，
\end{enumerate}

\[
\mu (f \cdot g ^ {1}) = \mu (f ^ {\chi} \cdot g).\tag{2}
\]

\begin{enumerate}
\def\labelenumi{\alph{enumi})}
\setcounter{enumi}{3}
\item
  若 \(f \in \mathcal{K}(\mathbf{X})\) 满足 \(f^{\chi} = 0\)，则 \(\mu(f) = 0\)。
\item
  \(\Rightarrow\) b)：若 \(\mu(f)=\mathrm{I}(f^{\chi})\)，则根据命题 1 d)，
\end{enumerate}

\[
\begin{array}{r l} & {\langle \pmb {\delta} (\xi) \mu , f \rangle = \langle \mu , \pmb {\delta} (\xi^ {- 1}) f \rangle = \mathrm{I} \Big (\big (\pmb {\delta} (\xi^ {- 1}) f \big) ^ {\chi} \Big)} \\ & {\qquad = \mathrm{I} \big (\chi (\xi) ^ {- 1} \Delta_ {\mathrm{H}} (\xi) f ^ {\chi} \big)} \\ & {\qquad = \chi (\xi) ^ {- 1} \Delta_ {\mathrm{H}} (\xi) \langle \mu , f \rangle ,} \end{array}
\]

从而 \(\delta (\xi)\mu = \chi (\xi)^{-1}\Delta_{\mathrm{H}}(\xi)\mu\)

\begin{enumerate}
\def\labelenumi{\alph{enumi})}
\setcounter{enumi}{1}
\tightlist
\item
  \(\Rightarrow\) c)：假设 b) 成立。注意，函数 \((x,\xi)\mapsto f(x)g(x\xi)\) 和 \((x,\xi)\mapsto f(x\xi)g(x)\) 定义在 \(\mathbf{X}\times\mathbf{H}\) 上，连续且支集紧（因为 H 在 X 上适当地作用）；据此，第三章第 4 节第 1 号定理 2 允许我们写出：
\end{enumerate}

\[
\begin{array}{r l} & {\int_ {\mathrm{X}} f (x) d \mu (x) \int_ {\mathrm{H}} g (x \xi) d \beta (\xi) = \int_ {\mathrm{H}} d \beta (\xi) \int_ {\mathrm{X}} f (x) g (x \xi) d \mu (x)} \\ & {\qquad = \int_ {\mathrm{H}} d \beta (\xi) \int_ {\mathrm{X}} f (x \xi^ {- 1}) g (x) \chi (\xi) \Delta_ {\mathrm{H}} (\xi) ^ {- 1} d \mu (x)} \\ & {\qquad = \int_ {\mathrm{X}} g (x) d \mu (x) \int_ {\mathrm{H}} f (x \xi^ {- 1}) \chi (\xi) \Delta_ {\mathrm{H}} (\xi) ^ {- 1} d \beta (\xi)} \\ & {\qquad = \int_ {\mathrm{X}} g (x) d \mu (x) \int_ {\mathrm{H}} f (x \xi) \chi (\xi) ^ {- 1} d \beta (\xi),} \end{array}
\]

这证明了 c)。

\begin{enumerate}
\def\labelenumi{\alph{enumi})}
\setcounter{enumi}{2}
\item
  \(\Rightarrow\) d)：若 c) 成立且 \(f^{\chi} = 0\)，则 \(\mu(f \cdot g^{1}) = 0\) 对所有 \(g \in \mathcal{K}(X)\) 成立，从而 \(\mu(f) = 0\)，只需选取 \(g \in \mathcal{K}(X)\) 使 \(g^{1} = 1\) 在 \(\operatorname{Supp} f\) 上成立（这可由将命题 2 应用于 \(\chi = 1\) 得到）。
\item
  \(\Rightarrow\) a)：若条件 d) 成立，则存在定义在 \(\mathcal{K}^{\chi}(\mathbf{X})\) 上的线性形式 I，使得 \(\mu(f) = \mathrm{I}(f^{\chi})\) 对 \(f \in \mathcal{K}(\mathbf{X})\) 成立，且根据命题 2，该形式相对有界。
\end{enumerate}

